\documentclass[11pt,a4paper]{article}
\usepackage[margin=22mm]{geometry}
\usepackage{amsmath,amssymb,booktabs,array}
\usepackage[hidelinks]{hyperref}
\setlength{\parindent}{0pt}
\setlength{\parskip}{5pt}
\raggedright
\emergencystretch=3em
\begin{document}
\begin{center}
{\Large Actual transfers on the restriction's own field}\\[4pt]
8 October 2026 --- analytical continuation and source-time decision
\end{center}
\textbf{Conclusion.} Specified one-interval fields close the traffic/receiving/
transfer/motion/count loop, including a local ramp extension. A nonuniform
donor exposes two different temporal interpretations of sampled requests.
No full multi-port numerical construction or useful-gradient result is claimed.

\section*{1. Common units, field and sampled requests}
Distance is metres, time seconds, and population continuous vehicle-equivalent
mass. Density $\rho$ has units eq/m; flow $F$ has units eq/s. Use
\begin{align}
 F(\rho)&=\min\{u\rho,w(\kappa-\rho)\},&
 u&=30,& w&=5,& \kappa&=3/20,\\
 \rho_c&=\frac{w\kappa}{u+w}=3/140,&
 Q&=u\rho_c=9/14,& L&=30,& T&=1.
\end{align}
Each section obeys $\partial_t\rho+\partial_xF(\rho)=0$ away from the
fixed ports and moving restrictions. Source IT moves from $(i,\ell)$ to
$(i+1,b)$ once. $\bar\rho_a=n_a/L$ is the donor section's average density;
$n_a$ is its mass. Native sampled sending/oriented requests are
\begin{align}
 S_a&=T\min(u\bar\rho_a,Q)=\min(n_a,QT),\\
 \theta_{ab}&=\pi_{ab}T\ (b\ne a),\quad
 \theta_{aa}=1-\sum_{b\ne a}\theta_{ab},\qquad d_{ab}=\theta_{ab}S_a.
\end{align}
$d_{ab}$ is desired interval mass, whereas $r_{ab}=d_{ab}/T$ is its nominal
rate. The source choice rate is $\pi_{ab}=(\bar v_b-\bar v_a)_+/(u\tau)$;
$\bar v$ denotes the approved past-period lookahead Edie observation, not
present point FD speed. Hold these choice observations and birth speeds at
the public clock. Numerical internal events do not introduce new births or
positive recovery. All records carry zero population.

For a fixed recipient port, its own right trace gives supply
$R(t)=\min\{w(\kappa-\rho(t,0+)),Q\}$. Donor own traces bound physically
available sending. Old-record actual speeds use their own one-sided ahead
traffic and the K.A running-minimum cap. In the examples below, that cap never
falls below the held desired speeds. Both candidate source-time interpretations
agree in these examples because donors stay available at the relevant rates.

Let $C(t)=n(0)+\int_0^t q_{\rm in}(s)ds$ be the local count gauge and define
\begin{equation}
 N(t,x)=C(t)-\int_0^x\rho(t,y)dy.
\end{equation}
Then $N(t,L)$ is cumulative local exit, $N_x=-\rho$, $N_t=F$, and along
an old record $z(t)$, $dN(t,z(t))/dt=F-\dot z\rho=0$ under zero relative
passing. These are the counts used below, not independently assigned labels.

\newpage
\section*{2. Constrained through/lateral entry: exact whole interval}
Initially the target section has $\rho_c$ throughout. An old zero-passing
record starts at $z(0)=1$ with held desired speed $s=5$. Own ahead traffic
is free or vacuum, so $z(t)=1+5t$. Its congested upstream state and flow are
\begin{equation}
 \rho_U=\frac{w\kappa}{w+s}=3/40,\qquad q_U=s\rho_U=3/8.
\end{equation}
The ordinary shock from $\rho_c$ to $\rho_U$ has speed
\begin{equation}
 \sigma=\frac{q_U-Q}{\rho_U-\rho_c}=-5,\quad
 b(t)=1-5t,\quad t_b=1/5.
\end{equation}
The downstream vacuum front is $e(t)=1+30t$, reaching $L$ at $t_e=29/30$.
For $0<t<t_b$, the own field consists of $\rho_c$ on $[0,b(t)]$,
$\rho_U$ on $[b(t),z(t)]$, vacuum on $[z(t),e(t)]$, and $\rho_c$ from
$e(t)$ to the exit when $e(t)<L$. For $t\ge t_b$, the queue starts at zero.
Thus the derived boundary fluxes, not externally supplied histories, are
\begin{equation}
 R(t)=q_{\rm in}(t)=
 \begin{cases}9/14,&0<t<1/5,\\3/8,&1/5<t<1,\end{cases}
 \quad
 q_{\rm out}(t)=\begin{cases}9/14,&0<t<29/30,\\0,&29/30<t<1.\end{cases}
\end{equation}
The two sampled donor demands are through $r_1=Q$ and lateral $r_2=Q/3$.
The lateral donor has choice fraction $1/3$ toward this target and $2/3$
toward its own next lane. For example $\tau=2$, Edie target observation 30
and donor observation 10 give this fraction. The fast donor's sampled staying
fraction is one. These are supplied history inputs; their reached-policy
origin is not asserted. Both donors start congested, with own sending $Q$.
IT gives
\begin{equation}
 f_1(t)=\tfrac34R(t),\quad f_2(t)=\tfrac14R(t),\quad
 f_{2\to\mathrm{own}}=\tfrac23Q=3/7.
\end{equation}
The parallel recipient starts free at density $1/70$ with free exit, so its
inflow and outflow are both $3/7$. No second advance is counted. Integrated
accepted masses and the target state are
\begin{align}
 I&=\int_0^1R(t)dt=3/7,& I_1&=9/28,& I_2&=3/28,& O&=87/140,\\
 n(1)&=9/14+3/7-87/140=9/20,\quad C(1)=15/14,\\
 z(1)&=6,\quad N(1,6)=15/14-9/20=87/140=N(0,1).
\end{align}
The ending field is $3/40$ on $[0,6]$ and vacuum ahead. The native coarse
step with initial receiving $Q$ would transfer $Q$ mass; the own causal
field transfers $3/7$. This is the declared receiving adaptation being
qualified, not a claim of identical coarse-step trajectories.

\newpage
\section*{3. Donor fields and a fully available receiving case}
To verify physical donor availability in section 2, start each 30 m donor
at density $1/10$, with inlet flow $1/4=F(1/10)$. Their downstream flows
are $f_1(t)$ and $f_2(t)+3/7$. Every flow is below $Q$. The congested
port density is $\rho^*(f)=\kappa-f/w$. On the congested branch all
resulting waves travel upstream at $-5$. At $t=1$ each donor has initial
density on $[0,25]$, first-period port density on $[25,26]$, and final-period
port density on $[26,30]$:
\begin{center}
\begin{tabular}{lccc}
\toprule
Donor & First port density & Final port density & Ending mass\\
\midrule
Through & $3/56$ & $3/32$ & $41/14$\\
Lateral + own through & $9/280$ & $51/1120$ & $19/7$\\
\bottomrule
\end{tabular}
\end{center}
All port densities are at least $\rho_c$ and below $\kappa$, hence sending
stays $Q$. Wave positions stay away from the inlet. The masses also equal
$3+1/4-I_a$, with $I_a=9/28$ and $15/28$. No unavailable mass is debited.

\textbf{Available case.} Use through/lateral rates $3/20$ and $1/10$; total
$q_0=1/4$. The target starts uniformly at $\rho_0=q_0/u=1/120$, with the
same old record at 1 m and desired speed 5. Its queue rear moves at
\begin{equation}
 \sigma_0=\frac{3/8-1/4}{3/40-1/120}=15/8>0.
\end{equation}
Consequently the inlet stays free: $R(t)=Q$, $q_{\rm in}=1/4$, both
requests are fully served. At the end its field is $1/120$ on $[0,23/8]$,
$3/40$ on $[23/8,6]$, and vacuum ahead; the original downstream free tail
has exited. Direct integration and conservation agree:
\begin{align}
 I&=1/4,\quad O=(1/4)(29/30)=29/120,\\
 n(1)&=(1/120)(23/8)+(3/40)(25/8)=31/120,\\
 C(1)&=1/2,\qquad N(1,6)=29/120=N(0,1).
\end{align}
The through donor is uniformly free at density $1/200$ and supplies $3/20$.
The lateral donor has saturated demand and choice fraction $7/45$, yielding
lateral $1/10$ and own through $19/35$. With the same $\tau=2$, target
Edie observation 30 and donor observation $62/3$ produce $7/45$. Its total
outflow is $Q$; a congested initial field can sustain that with upstream
waves. Its parallel recipient starts free at $(19/35)/30=19/1050$ and
receives/exits $19/35$. Both cases honor the separate past-speed convention.

\newpage
\section*{4. Local ramp, multiple records and sampled births}
Return to section 2's target. Add a saturated finite ramp. Let its fixed
priority be illustratively $p=1/3$, unchanged within the calculation.
Non-ramp total desired rate is $4Q/3$; it includes both through and lateral
traffic. Ramp demand is $Q$. The approved shared node gives
\begin{equation}
 f_r=R/3,\quad f_m=2R/3,\quad
 f_{\rm through}=R/2,\quad f_{\rm lateral}=R/6.
\end{equation}
All three still total $R$; receiving is spent once. Integrated masses are
$1/7,3/14,1/14$. With 30 m congested donors initially at $1/10$ and
inlets $1/4$, all resulting port states are congested and admissible.
Their ending masses are, respectively, $87/28,85/28,11/4$; the lateral
donor's $11/4$ also debits its parallel through transfer $3/7$.
The ramp's upstream receiving stays $1/4$ because its downstream waves
cannot reach the upstream boundary in one second. M.A may supply $1/4$
there under adequate offered flow and permission. This is compatible upstream
admission, not a new contemporaneous cap on actual merge entry.

\textbf{A second old record.} Put another zero-passing record at 2 m, with
held desired speed 15. Its own upstream state is $\rho_{U2}=3/80$.
Initially its queue rear is $2-5t$, front $z_2(t)=2+15t$, downstream vacuum
front $2+30t$. The first record's vacuum front meets this queue rear at
\begin{equation}
 1+30t=2-5t\quad\Longrightarrow\quad t_*=1/35.
\end{equation}
At that contact, the ordinary vacuum-to-$\rho_{U2}$ shock has speed 15;
its continuing rear is $b_2(t)=10/7+15t$. Its queue therefore has constant
length $4/7$ and mass $3/140$ after $t_*$. No passing or coincident-record
jitter is introduced. The first record stays at speed 5; any encountered
second queue moves at 15 and does not require slower braking. The second's
ahead traffic is free/vacuum, so its own cap permits 15. The inlet wave in
section 2 is unaffected. Exit flow is $Q$ until $14/15$, then zero.
At the end the target field is
\begin{equation}
 \rho(1,x)=\begin{cases}
 3/40,&0<x<6,\\3/80,&115/7<x<17,\\0,&\text{elsewhere in }(0,30),
 \end{cases}
 \quad n(1)=9/20+3/140=33/70.
\end{equation}
Conservation gives $9/14+3/7-3/5=33/70$. The count labels are
\begin{equation}
 C(1)=15/14,\quad h_1=N(1,6)=87/140=N(0,1),\quad
 h_2=N(1,17)=3/5=N(0,2).
\end{equation}
All queue and gap masses agree with these counts.

\textbf{Completed E1 attachment.} For illustrative threshold one equivalent,
past ramp/lateral counters $9/10$ and $19/20$ become $73/70$ and $143/140$.
Each crosses one threshold; remainders $3/70$ and $3/140$ stay recorded.
At $t=1$ insert two distinct records at the recipient inlet, both with count
$N(1,0)=15/14$, retaining request-phase desired speeds (e.g. 5 and 10).
They add no mass and do not change preceding receiving. The own queue ahead
has speed 5, so its trace supplies the immediate next-interval cap. Neither
arbitrary later contact continuation nor section crossings are proved here.

\newpage
\section*{5. Nonuniform sending: a causal bound and two distinct completions}
Take a finite ramp approach of length 30 with initial $\rho=3/20$ on
$[0,3]$ and vacuum on $(3,30)$. Close upstream admission this interval.
Give its merge ample receiving, with no competing mainline demand; N.A then
reallocates the unused mainline share. Initial mass is $n_0=9/20$ and
the source average gives
\begin{equation}
 S=T\min(u n_0/L,Q)=9/20.
\end{equation}
The ordinary jam-to-vacuum triangular fan has edges $3-5t$ and $3+30t$,
with critical density between. The first traffic reaches the exit at
\begin{equation}
 t_a=(30-3)/30=9/10,\qquad
 \int_0^1 f(t)dt\le Q(1-t_a)=9/140<9/20.
\end{equation}
This bound follows from the own field; an interpolation cannot physically
deliver the full coarse accepted $S$ in this interval. CTM's original uniform
internal approximation is not refuted by this comparison.

\textbf{T.A: interval mass ceiling.} Let actual flow be the available own
flow until total request $S$ is spent, also respecting receiving. Here
$f_A(t)=0$ before $t_a$ and $Q$ after; the ceiling never binds. Thus
\begin{equation}
 I_A=9/140,\quad n_A(1)=9/20-9/140=27/70.
\end{equation}
With the closed inlet, the rear fan reaches $x=0$ at $3/5$ s. The ensuing
vacuum-to-critical rear moves at $u$, reaching $x=12$ at $t=1$. The ending
field is $\rho_c$ on $[12,30]$, vacuum before. Its mass is $(18)(3/140)=27/70$.

\textbf{T.B: held nominal request rate.} Impose additionally $f_B\le S/T$,
while never exceeding own availability. Here $f_B=0$ until $t_a$, then
$9/20$ instead of $Q$. Its own outlet queue state is
\begin{align}
 \rho_B^*&=\kappa-(S/T)/w=3/50,\quad
 I_B=(9/20)(1/10)=9/200,\\
 n_B(1)&=9/20-9/200=81/200.
\end{align}
The queue rear travels at $-5$ from the exit after $t_a$; it reaches
$x=59/2$ at $t=1$. The ending field is $\rho_c$ on $[12,59/2]$ and
$3/50$ on $[59/2,30]$, with vacuum before. Its mass is exactly $81/200$.
Both completions are conservative, causal and jam-bounded; the second has
an additional nominal-rate constraint. They differ by $27/1400$ mass.
An initially empty 30 m recipient with free exit contains the transferred
traffic on $[0,3]$ at the ending clock: density $3/140$ under T.A or $3/200$
under T.B. Its mass is respectively $I_A$ or $I_B$, with no exit yet. Thus
these are paired completed merge entries, not merely donor debits.

For an eligible actual ramp-entry counter initially $H_0=47/50$, E1 yields
\begin{equation}
 H_A(1)=703/700>1,\qquad H_B(1)=197/200<1.
\end{equation}
T.A creates one next-clock effect (remainder $3/700$), whereas T.B creates
none. This is an operator diagnostic of the approved birth threshold, not
an optimized-control witness or a proposal to manufacture hidden information.

\newpage
\section*{6. Why the mass ceiling is natural, and what it does not settle}
For a section with $L=uT$, all upstream-port admissions strictly after the
current clock need at least $T$ to traverse its full length because actual
forward speeds do not exceed $u$. They cannot leave this section within the
same regular interval. Lane transfers also enter its upstream port. With
conservation and ordinary stationary-port capacity $Q$,
\begin{equation}
 \text{actual total outgoing mass}\le \min\{n(0),QT\}=S.
\end{equation}
This excludes interior sources and zero-length shortcuts, neither of which
is approved here. It supports T.A's total sending ceiling without requiring
a second held rate. It does not assign exact movement-specific allowances
under selective target blockage: $d_{ab}=\theta_{ab}S_a$ still needs
explicit unspent-claim accounting, donor joint availability and target
competition. A naive repeated use of instantaneous choice fractions may
retry the same traffic and exceed a movement's sampled desired mass.

For T.A, proposed numerical claim state is
\begin{equation}
 b_{ab}(t)=d_{ab}-\int_0^t f_{ab}(s)ds\ge0.
\end{equation}
It is not extra population or a vehicle identity. The common constraints
include $\sum_b f_{ab}\le D_a^{\rm own}$, target shared receiving, and zero
flow when its claim is exhausted. N.A and IT competition must have a declared
meaning while positive claims are temporarily unavailable or become depleted.
Those constraints alone do not define the full multi-port allocation. T.B
instead holds $r_{ab}=d_{ab}/T$ as nominal rates, additionally reducing them
for unavailable donor traffic or constrained receiving; it likewise needs a
qualified shared-port transcription. No new upstream control is proposed.

The methods differ at finite public resolution. Their relation to the
source's approximation/refinement and their coupled motion/birth histories
must be qualified, not asserted from dimensional conversion. Because they
change those histories, Q.J.F.T is a shared methodological choice rather
than a private solver setting. T.A is provisionally preferred, not adopted.

\section*{7. Primary basis and evidence limits}
Laval and Daganzo (2006), \emph{Lane-changing in traffic streams},
sections 2.1--2.2, equations 5--13 and footnote 2, supplies desired sampled
advancing masses, Edie observations and IT. Its original numerical/source
recipe is not claimed to prove the present adaptation.

Jin, arXiv:0810.3952v2, sections 3--5, distinguishes discrete local merge
fluxes and continuous Riemann solutions. Interior states can represent the
limit; section 4.1 proves the fixed-priority flux's invariance. A simple
two-incoming-road theorem does not establish the present oriented multilanes.
No new priority is inferred for IT from that theorem.

\href{https://doi.org/10.1016/j.trb.2005.04.003}{Lane-changing journal DOI}
\quad
\href{https://arxiv.org/pdf/0810.3952}{Continuous merge primary manuscript}

The examples use admissible specified initial/history state; no reached
public scenario, arbitrary contact/crossing solver, clearance, full TST or
useful gradient is inferred. The four-step pass stops at Q.J.F.T. A complete
qualified finite construction can still meet Q.J.F without a universal
continuum existence theorem. Existing M.A/N.A/G1/E1/K.A remain approved;
representation A1/A2/C and the complete law remain unadopted. Stage 4 unopened.
\end{document}
